
\section{CORD-19}\fbeg{}\begin{frame}{\ft{Case Study: \q{CORD-19}}}

\begin{ftxt}[.9\textwidth]

The CORD-19 corpus, developed by the Allen Institute for 
Artificial Intelligence (AI2), is a good example of the 
goals and techniques for text-mining in the context of 
scientific research.  This corpus included over 400,000 
freely-available full-text articles available both 
in human-readable formats and in machine-readable 
encodings --- specifically, S2ORC JSON 
(S2ORC comes from Semantic Scholar Open Research Corpus).  
Because Covid research spanned many disparate subject 
areas --- viral morphology, infectivity, genetic profiles 
and strains, vaccine development, diagnostics, clinical 
treatments, epidemiology, etc. --- automated text-mining 
tools could, at least in principle, help scientists 
discover research in other disciplines to help 
fit the pandemic's pieces together.  AI2 did not 
specifically create those tools, but they curated 
CORD-19 to be directly accessable from computer 
software in the hopes that other parties could 
build text-mining capabilities upon that basis.

\end{ftxt}

\end{frame}

